\chapter{Plano carlesino}

\section{Canales direccionales}
\begin{definition}
El plano carlesino se fija como una estructura de seis canales direccionales
\[
(+x,-x,+y,-y,+z,-z),
\]
tratados como componentes independientes.
\end{definition}

\begin{observation}
La direccionalidad no se reduce aquí a un signo sobre un único eje. Cada par de opuestos se declara como doble canal estructuralmente distinguible.
\end{observation}

\section{Estado direccional}
\begin{definition}
Un estado direccional se representa mediante el 6-vector
\[
p=(p_{+x},p_{-x},p_{+y},p_{-y},p_{+z},p_{-z}).
\]
\end{definition}

\begin{definition}
El origen canónico se identifica con el estado nulo
\[
\subz \equiv (0,0,0,0,0,0).
\]
\end{definition}

\begin{observation}
Esta representación fija una lectura primitiva por canales. No impone todavía una métrica completa ni una interpretación física del espacio.
\end{observation}

\section{Involución estructural de opuestos}
\begin{definition}
La involución mínima actúa sobre el estado direccional intercambiando canales opuestos:
\[
\Iop(p_{+x},p_{-x},p_{+y},p_{-y},p_{+z},p_{-z})
=
(p_{-x},p_{+x},p_{-y},p_{+y},p_{-z},p_{+z}).
\]
\end{definition}

\begin{observation}
Aplicada dos veces, la involución recupera el estado inicial. En este bloque, su función es declarar simetría estructural entre opuestos dentro de la representación por canales.
\end{observation}

\section{Lecturas compatibles}
\begin{definition}
La estructura de seis canales admite dos niveles de lectura compatibles:
\begin{enumerate}
  \item una lectura primitiva por canales no negativos;
  \item una lectura proyectada derivada, obtenida por diferencia entre canales opuestos.
\end{enumerate}
\end{definition}

\begin{observation}
La lectura proyectada no reemplaza la lectura primitiva. Ambas capas se consideran compatibles siempre que la proyección se declare como derivada y no como estructura originaria.
\end{observation}

\section{Cláusula de alcance}
\begin{observation}
La introducción del plano carlesino fija una organización mínima de direcciones, oposición y estado direccional sin colapsar memoria de orientación. No valida por sí sola toda geometría ulterior del sistema.
\end{observation}

